\SourcePage{325}{328}
\section{Wave functions of outer electrons near the nucleus}

We have seen (from the Thomas--Fermi model) that the outer electrons of complex atoms (large $Z$) are found mainly at distances $r\sim1$ from the nucleus\footnote{Atomic units are used in this section.}. A number of atomic properties, however, depend essentially on the electron density close to the nucleus (we shall encounter such properties in \S\S~72 and 120). To determine the order of magnitude of this density, let us trace how the electron wave function $\psi(r)$ in an atom changes as $r$ decreases from large distances ($r\sim1$) to small ones.

In the region $r\sim1$, the nuclear field is screened by the other electrons, so that the potential energy $U(r)\sim1/r\sim1$. The energy $E$ of an electron level in this field is also of order unity. At distances of the order of the Bohr radius in a field of charge $Z$ ($r\sim1/Z$), on the other hand, the nuclear field may be regarded as unscreened: $U=-Z/r$. In the intermediate region, $1/Z\ll r\ll1$, the potential energy $|U|$ is already large compared with the electron energy $E$, and the condition
\[
 \frac{d}{dr}\frac1p\sim\frac{d}{dr}\frac1{\sqrt{|U|}}\ll1
\]
is satisfied ($p$ is the momentum), so the motion is semiclassical. The spherically symmetric semiclassical wave function is
\begin{equation}
 |\psi(r)|\sim\frac1{r\sqrt p}\sim\frac1{r|U|^{1/4}},\qquad \frac1Z\ll r\ll1,
 \tag{71.1}
\end{equation}
and the order of magnitude of its coefficient ($\sim1$) is fixed by matching to the wave function, $\psi\sim1$, at $r\sim1$.

\SourcePage{326}{329}
Using (71.1) as an order-of-magnitude expression at $r\sim1/Z$ (with $U=-Z/r$ substituted into it), we obtain the required value of the wave function near the nucleus\footnote{To determine the coefficient in this formula (given the wave function in the region $r\sim1$), expression (36.25) would have to be used in the region $r\lesssim1/Z$.}:
\begin{equation}
 \psi\left(\frac1Z\right)\sim\sqrt Z.
 \tag{71.2}
\end{equation}
In accordance with the general properties of wave functions in a central field (\S~32), as the distance is reduced further, $\psi(r)$ either remains constant in order of magnitude (for an $s$ electron) or begins to decrease (when $l\ne0$).

The probability of finding the electron in the region $r\lesssim1/Z$ is
\begin{equation}
 w\sim|\psi|^2r^3\sim\frac1{Z^2}.
 \tag{71.3}
\end{equation}
The formulas (71.2) and (71.3), of course, specify only the systematic trend of these quantities as $Z$ increases, without allowing for nonsystematic changes from one element to the next.
